% !TeX root = ./Thesis.tex

% 第五章
\chapter{原型系统仿真评估结果与分析}
\label{chapter:chapter05}

\section{仿真运行概况与稳定性分析}
\label{sec:5.1}

\subsection{测试用例与对话会话生成规模}
\label{subsec:5.1.1}

为了全面评估基于大语言模型的初中物理迷思概念干预系统的教学有效性与护栏约束效果，本研究开展了规模化的系统仿真实验。本章所分析的数据均来自一次真实的系统测试套件运行（仿真套件目录：\texttt{suite\_2026-04-16\_14-42-05}），该套件包含了五个独立重复的运行批次（\texttt{run\_01} 至 \texttt{run\_05}），以保证结果的统计稳定性。

在测试用例设计上，仿真实验覆盖了第四章所述的初中物理两类核心迷思主题：电学（如“序列衰减观” M-ELE-001 和“回路消耗观” M-ELE-002）与浮力（如“重物必沉单因归因” M-BUO-001 和“浮力等于深度压强经验投射” M-BUO-002）。同时，实验引入了三种典型的模拟学生画像（P1：固执型/逻辑熔断易发者；P2：易受经验误导型；P3：高认知负荷/易挫败型），以此来模拟真实教学场景中学生认知状态的多样性。

实验共对比了系统的三个版本：
\begin{itemize}
    \item \textbf{Baseline}：未引入第四章所述阶段化控制与护栏机制的基础大语言模型（仅使用基础系统提示词）。
    \item \textbf{FSM}：引入了阶段化教学主线（T1--T5，对应 \texttt{S0-S7} 状态流转）的版本，但不包含安全护栏与支架豁免机制。
    \item \textbf{FSM+Guardrail}：第四章所设计的完整原型系统，在阶段化教学主线的基础上，融合了输入侧拦截、输出侧认知卸载审查以及认知死锁降级干预机制（即护栏分支 G1--G4 及 \texttt{ANTI\_LOOP\_RULES}）。
\end{itemize}

每个运行批次中，每种系统版本均与所有（学生画像 $\times$ 迷思主题）组合进行完整会话，各版本在单次 Run 中生成一定数量的对话轮次。汇总五个运行批次的结果，实验共收集了大量带有精细状态标签（如诊断分类置信度、护栏触发次数、答案泄露次数等）的会话记录。这些记录为后续的微观干预轨迹分析和宏观统计比较提供了坚实的数据基础。

\subsection{系统运行层指标结果（异常终止率、响应稳定性等）}
\label{subsec:5.1.2}

在深入分析教学效果之前，必须首先验证系统作为教育智能体的运行稳定性。表~\ref{tab:system-stability} 报告了三个系统版本在异常终止率（Abnormal Termination Rate）、平均对话轮次（Average Turns）以及状态转移成功率（Transition Success Rate）等维度的总体表现。数据取自五个实验批次的均值（附 95\% 置信区间）。

\begin{table}[htbp]
    \centering
    \caption{各版本系统运行层稳定性指标对比（$N=5$）}
    \label{tab:system-stability}
    \begin{tabular}{lccc}
        \hline
        \textbf{指标} & \textbf{Baseline} & \textbf{FSM} & \textbf{FSM+Guardrail} \\
        \hline
        异常终止率 & $0.00\% \pm 0.00\%$ & $0.00\% \pm 0.00\%$ & $0.00\% \pm 0.00\%$ \\
        平均对话轮次 & $9.95 \pm 0.54$ & $9.15 \pm 0.55$ & $9.03 \pm 0.20$ \\
        状态转移成功率 & $13.76\% \pm 1.91\%$ & $30.19\% \pm 2.07\%$ & $31.10\% \pm 1.66\%$ \\
        \hline
    \end{tabular}
\end{table}

从表~\ref{tab:system-stability} 中可以看出，所有版本的异常终止率均为 0.00\%，表明底层的 API 调用和基本对话循环是高度稳健的，未出现因系统错误导致对话意外中断的情况。

在平均对话轮次方面，Baseline 版本的平均轮次最长（$9.95$ 轮），而引入阶段化控制后（FSM 和 FSM+Guardrail），平均对话轮次有所下降（分别为 $9.15$ 轮和 $9.03$ 轮），且 FSM+Guardrail 版本的标准差与置信区间显著缩小（$\pm 0.20$）。这说明阶段化教学主线有效地收束了对话的发散性，使得干预过程更加紧凑、可控；护栏机制的加入进一步避免了冗长且无意义的相互绕圈子。

最显著的差异体现在状态转移成功率上。Baseline 版本的状态转移成功率仅为 $13.76\%$，这说明无结构控制的大模型往往难以自觉地遵循“诊断$\rightarrow$冲突$\rightarrow$重构$\rightarrow$验证”的教学递进逻辑，容易在某一状态中停滞。而 FSM 版本的成功率提升至 $30.19\%$，FSM+Guardrail 版本进一步提升至 $31.10\%$。这有力地验证了第四章中“阶段化教学对话控制（S0--S8 状态机）”设计的有效性，系统能够更顺畅地牵引学生从迷思暴露向概念验证推进。

\section{阶段化对话控制对干预引导的影响}
\label{sec:5.2}

\subsection{迷思概念识别准确率对比（Baseline vs. 阶段化控制）}
\label{subsec:5.2.1}

迷思概念的精准识别是后续发起针对性认知冲突的基础。系统不仅需要识别出学生回答错误，还需要准确分类其错误的深层原因（如识别出其属于“序列衰减观”而非“闭合回路混淆”）。表~\ref{tab:identification-accuracy} 展示了不同系统版本在迷思概念识别准确率（Identification Accuracy）上的对比。

\begin{table}[htbp]
    \centering
    \caption{迷思概念识别准确率对比}
    \label{tab:identification-accuracy}
    \begin{tabular}{lccc}
        \hline
        \textbf{版本} & \textbf{识别准确率均值} & \textbf{标准差} & \textbf{95\% 置信区间} \\
        \hline
        Baseline & 96.84\% & 2.15\% & $\pm 1.88\%$ \\
        FSM & 96.22\% & 1.88\% & $\pm 1.65\%$ \\
        FSM+Guardrail & 91.49\% & 7.64\% & $\pm 6.70\%$ \\
        \hline
    \end{tabular}
\end{table}

数据显示，无论是 Baseline 还是引入状态机（FSM）的版本，大语言模型在迷思概念识别上均表现出极高的基础准确率（超过 96\%）。值得注意的是，FSM+Guardrail 版本的识别准确率略有下降（91.49\%）且方差增大。深入日志分析发现，这种现象并非由于模型理解能力退化，而是由于输入侧护栏（Guardrail）的严格拦截机制。当学生输入包含“直接索答”或情绪性抗拒时，护栏机制（G1--G2分支）会优先将其拦截并进行温和重定向（Refusal/Guidance），导致系统在这一轮次没有进行实质性的物理概念分类，从而在统计层面上拉低了单纯的“诊断准确率”。这一结果表明，教学安全与单纯的语义诊断之间存在一定的权衡（Trade-off）。

\subsection{认知修正率与教学轮次推进对比}
\label{subsec:5.2.2}

概念干预的最终目的是帮助学生放弃迷思概念并建立科学概念。实验使用“认知修正率（Cognitive Correction Rate）”作为衡量这一结果的核心指标。该指标反映了在对话结束时，学生能否在验证阶段（S6/S7）用自己的话正确表述物理原理，并在模拟后测中通过测试。

\begin{table}[htbp]
    \centering
    \caption{认知修正率与拒绝成功率对比}
    \label{tab:correction-rate}
    \begin{tabular}{lcccc}
        \hline
        \textbf{版本} & \textbf{认知修正率均值} & \textbf{95\% 置信区间} & \textbf{拒绝索答成功率} & \textbf{95\% 置信区间} \\
        \hline
        Baseline & 56.67\% & $\pm 8.00\%$ & 0.00\% & $\pm 0.00\%$ \\
        FSM & 60.00\% & $\pm 6.11\%$ & 0.00\% & $\pm 0.00\%$ \\
        FSM+Guardrail & 58.33\% & $\pm 7.31\%$ & 20.00\% & $\pm 39.20\%$ \\
        \hline
    \end{tabular}
\end{table}

如表~\ref{tab:correction-rate} 所示，Baseline 版本的认知修正率为 56.67\%，FSM 版本提升至 60.00\%。这一提升证实了第四章提出的假设：有组织的阶段化教学（从 T1 澄清观念到 T3 制造冲突，再到 T4 支架过渡）比无结构的自由问答更能有效促成概念转变。FSM 减少了漫无目的的讨论，使对话轮次（如~\ref{subsec:5.1.2} 节所示）更高效地服务于认知重构。

FSM+Guardrail 版本的修正率（58.33\%）介于两者之间。虽然略低于纯 FSM 版本，但该版本取得了 20.00\% 的拒绝索答成功率（Refusal Success Rate，而 Baseline 和纯 FSM 均为 0\%）。这意味着 FSM+Guardrail 牺牲了少量的“表面修正率”（部分固执型学生因未能直接获取答案而未完成最终修正），但换来了对“认知卸载”行为的有效防范，确保了最终通过验证的学生是真正经历了独立思考的。

\subsection{自动化模拟评估与人工抽核：教学有效性维度对比}
\label{subsec:5.2.3}

为进一步验证系统在微观互动中的教学有效性（Teaching Effectiveness），本研究从 \texttt{run\_01} 中随机抽取了多段对话记录（\texttt{manual\_audit.csv}）进行人工与大语言模型裁判（LLM-as-a-Judge）的双重评估。评估维度主要聚焦于“是否成功引导学生产生认知顿悟”。

在对 \texttt{sim\_Baseline\_P1\_M-ELE-001\_8db7b7} 的审核中，Baseline 系统（评分为 2 分）虽然频繁使用反问和类比（如水车、送货员等），但始终未能切中学生“电能与电荷流动”混淆的要害。由于缺乏阶段控制，系统在学生仍坚持错误观念时便陷入了比喻的堆砌，最终未能帮助 P1（固执型）学生突破认知僵局。

相比之下，在 \texttt{sim\_FSM+Guardrail\_P2\_M-BUO-001\_809f87}（浮力重物必沉迷思）中，系统精准地执行了阶段化控制：首先在 S3 阶段诊断出单因归因错误，随后在 S4 阶段抛出极端反例（“钢铁巨轮为何能浮，小铁钉为何沉底？”）引发认知冲突。在学生产生困惑（“是不是船里面是空的”）后，系统没有直接告知阿基米德原理，而是提供了一个低负荷的思想实验支架（“如果把巨轮揉成实心铁球会怎样？”），成功引导学生认识到“排开液体体积”才是关键。这种步步为营的引导，在人工抽核中获得了较高的教学有效性评分，充分体现了阶段化控制在复杂概念教学中的优势。

\section{教学安全约束机制的效果分析}
\label{sec:5.3}

\subsection{答案泄露率（Answer Leakage Rate）对比}
\label{subsec:5.3.1}

生成式人工智能在教育应用中的核心风险之一是沦为“作业代写机”，直接向学生泄露最终答案，从而剥夺学生的认知劳动。本研究在第四章中设计了输出侧认知卸载审查机制。表~\ref{tab:leakage-rate} 呈现了该机制拦截答案泄露的直接效果。

\begin{table}[htbp]
    \centering
    \caption{各版本系统答案泄露率对比}
    \label{tab:leakage-rate}
    \begin{tabular}{lccc}
        \hline
        \textbf{版本} & \textbf{答案泄露率均值} & \textbf{标准差} & \textbf{95\% 置信区间} \\
        \hline
        Baseline & 7.13\% & 4.29\% & $\pm 3.76\%$ \\
        FSM & 8.94\% & 1.77\% & $\pm 1.55\%$ \\
        FSM+Guardrail & 0.00\% & 0.00\% & $\pm 0.00\%$ \\
        \hline
    \end{tabular}
\end{table}

数据清晰地表明，未加护栏约束的 Baseline 和单纯状态机控制的 FSM 版本，其答案泄露率分别达到了 7.13\% 和 8.94\%。在这些版本中，当学生反复表达困惑或直接索要公式时，大语言模型往往会出于其内置的“帮助人类用户完成任务”的对齐倾向，直接输出完整的推导过程或最终结论。

而引入安全约束机制的 FSM+Guardrail 版本，其答案泄露率被完美控制在 0.00\%。这一结果强有力地证明了我们在第四章中设计的输出审查漏斗（基于 LLM-as-a-Judge 的 Leakage Evaluation）能够严格守住“不直接给答”的底线。

\subsection{拦截触发特征与启发式豁免效果分析}
\label{subsec:5.3.2}

“零泄露”并非通过简单粗暴的“拒绝回答”来实现的。根据聚合数据，FSM+Guardrail 版本的平均护栏拦截触发率（Guardrail Interception Rate）为 2.60\%（$\pm 1.35\%$）。这意味着系统并非时刻在拦截，而是在关键节点上发挥了兜底作用。

更重要的是，第四章中提出的“教学支架豁免”机制在保持低泄露率的同时，维系了对话的流畅度。在 \texttt{sim\_FSM+Guardrail\_P3\_M-BUO-002\_aade88}（压强经验投射迷思）日志中，面对 P3（易挫败型）学生在上下表面压力差问题上的卡顿，系统成功触发了【解释类比/思想实验】的豁免：“假设原来上面受压 10N，下面受压 15N...现在都增加 100N...新的差值是多少？”。系统通过给出具体的计算支架，帮助学生跨越了抽象符号推演的鸿沟，而没有直接告诉其“浮力等于排开液体重力”的最终结论。这种在“安全边界内提供最小充分支架”的策略，有效缓解了认知过载，避免了对话因过度严格的拦截而陷入死锁。

\subsection{自动化模拟评估与人工抽核：苏格拉底式提问特征符合度对比}
\label{subsec:5.3.3}

在对微观干预日志的审计中（\texttt{manual\_audit.csv}），大语言模型裁判对所有版本的“苏格拉底度（Socratic Degree）”均给出了最高分（5分）。这表明无论是否有护栏，基础模型在提示词工程的引导下，都具备了采用反问、类比和启发式探究的基本语域能力。

然而，护栏版本的优越性体现在提问的“韧性”上。例如，在 \texttt{sim\_FSM\_P3\_M-ELE-001\_fd94e4} 案例中，纯 FSM 版本虽然也是以提问为主，但当对话陷入深水区时，出现了多达 4 次的轻微答案泄露。而在对应的 FSM+Guardrail 版本中，系统不仅保持了苏格拉底式的循循善诱，还在学生面临“顿悟”的临界点上，坚守住了底线，逼迫学生自行说出关键机制（如“原来只要托力够大，就算很重也能浮着啊”），从而实现了更高质量的认知参与。

\section{对话干预的差异表现与边界现象}
\label{sec:5.4}

\subsection{阶段化控制+护栏版本下的优质干预案例}
\label{subsec:5.4.1}

日志档案中的 \texttt{sim\_FSM+Guardrail\_P3\_M-ELE-002} 展现了一次教科书级别的电学迷思干预。学生 P3 持有典型的“单向电流/无需闭合回路”迷思（“电从正极射出来进到灯泡里就够了”）。

在干预过程中，系统首先在 T2 阶段（诊断）澄清了学生的错误表征。随后进入 T3 阶段（认知冲突），系统运用了\textbf{后果探索策略（归谬）}：“如果正极像机关枪一样发射电荷而负极不回收，电池里的电荷能无限发射吗？”当学生意识到电荷会耗尽时，系统进一步引入了\textbf{类比支架}：“就像排队进电影院，只进不出，门很快就会堵死”。最终，在 T4 和 T5 阶段，系统帮助学生明确区分了“电荷流动（快递员）”与“能量转化（包裹）”的核心差异。该案例中，系统既没有泄露答案（Leakage=0），也没有触发异常（Abnormal End=false），仅用 5 轮对话就引导学生完成了从“单向消耗”到“循环回路”的概念转变\cite{kucukozer2007misconceptions}，体现了系统在处理复杂抽象物理概念时的卓越能力。

\subsection{固执型学生画像（P1）下的引导失败或逻辑熔断案例分析}
\label{subsec:5.4.2}

尽管系统在多数情况下表现优异，但针对 P1（固执型/逻辑熔断易发者）的干预仍暴露出明显的边界挑战。数据表明，在多个针对 P1 画像的测试中（如 \texttt{sim\_FSM+Guardrail\_P1\_M-ELE-001\_c645fb} 和 \texttt{sim\_FSM+Guardrail\_P1\_M-BUO-002\_627f55}），对话轮次均达到了系统设定的上限（13轮），且最终状态停留在“认知僵局（Cognitive Deadlock）”，终止原因为 \texttt{max\_turns\_reached}。

人工抽核意见指出，在这些失败案例中，系统虽然“在苏格拉底度上表现极佳，通过连续反问（如串联100个灯泡、快递员比喻等）引导学生暴露认知矛盾，避免直接给出答案”，但未能有效回应 P1 学生特有的顽固核心困惑（如“能量没了电流为何不变”）。由于护栏机制严格禁止直接告知结论，系统在第四章设计的“认知死锁降级干预机制（\texttt{ANTI\_LOOP\_RULES}）”被触发后，虽然回退到了提供微支架的 S5 阶段，但如果微支架（如受控思想实验）仍无法击穿学生的认知壁垒，系统便会在 S5 与学生原地“拉锯”，最终耗尽轮次。

这一边界现象深刻揭示了完全基于文本对话的智能导师的局限性：对于具有强烈既有经验执念的学生，仅靠语言层面的逻辑归谬和隐喻类比可能不足以动摇其底层解释框架\cite{posner1982accommodation,chi1994things}。在未来的迭代中，此类情形可能需要转介至更直观的多模态交互（如提供可操作的虚拟仿真电路板），或触发更高权限的“直接讲授与事实覆盖”，以打破语言逻辑的死锁。这为教育智能体的发展指明了从单一对话模式向多模态具身认知支持演进的研究方向。
